\section*{Task 1: Fixed-arrow tangent sequence}

\paragraph{Goal.}
Record the deformation-theoretic exact sequence for the fixed-arrow slice and
clarify how it compares with the final codimension argument.

\paragraph{Deformation-theoretic picture.}
Let
\[
\Rep_d(A_L^{\mathrm{wt}})
\]
be the ambient free-arrow representation variety and
\[
\pi:\Rep_d(A_L^{\mathrm{wt}})\to \Hom(V_L,V_0),
\qquad
\phi\mapsto \phi_0,
\]
the projection to the closing-arrow component. The fixed-arrow slice is the
fiber
\[
\Rep_{d,W_0}(Q,I_L^{\mathrm{wt}})=\pi^{-1}(W_0),
\]
which is identified with the cyclic locus \(\Gamma_{d,S}\).

At a point \(M\) with closing arrow \(W_0\), the usual Voigt-type argument gives
an exact sequence
\[
0\to T_M\Gamma_{d,S}/T_M(G_d^{W_0}\cdot M)
\to \Ext_A^1(M,M)
\xrightarrow{\delta_{W_0}}
\Hom(V_L,V_0)\big/
T_{W_0}\bigl((\GL(V_0)\times \GL(V_L))\cdot W_0\bigr).
\]
If the boundary map \(\delta_{W_0}\) is surjective and the slice is smooth at
\(M\), this yields
\[
\operatorname{codim}_{\Gamma_{d,S}}(G_d^{W_0}\cdot M)
=
\dim \Ext_A^1(M,M)-(d_0-s)(d_L-s),
\]
where
\[
s=\operatorname{rank}(W_0)=d_L-r.
\]

\paragraph{Why the final theorem will not rely on this.}
The orbit decomposition proved earlier is already exact: every relevant stratum
is a single group orbit. Therefore the codimension of each stratum can be
computed directly from the dimension of the acting group and the dimension of
the stabilizer. This removes the need for smoothness or transversality
hypotheses. The tangent-sequence viewpoint is still useful as a consistency
check, but it is no longer the main engine of the codimension formula.

\paragraph{Status.}
Completed. The deformation-theoretic viewpoint is recorded, but the final proof
will use exact orbit-dimension formulas instead.
